% ---
% Abstract
% ---

\begin{resumo}[Abstract]
 \begin{otherlanguage*}{english}
% -----------------------------------------------------------------
% English version of the Resumo. Same structure and length.
% -----------------------------------------------------------------

Wire breakage in enameled copper and aluminum wire factories is an operational symptom whose origin may lie either in defects embedded in the raw material or in degraded equipment conditions, with fracture phenotypes often compatible with more than one hypothesis and without standardized cause labeling in the operational records. This dissertation presents the qualification stage of WireSense, a hybrid decision-support tool proposed to infer the likely origin of each breakage event by combining two independent strands of evidence. The method was developed on data from a factory operating with seventy-seven enameling machines in parallel and is structured into two pillars. The Operational Signals Pillar applies anomaly detection via \emph{Isolation Forest} over sensor time series, with one model trained per machine on the history of intervals without incidents. The Historical Statistics Pillar operates on the status table joined with the production orders, producing two comparative indicators per batch: the batch-relative rate, standardized by machine-hours accumulated per material, and the cross-machine failure score, which measures the spread of breakages across machines effectively exposed to the batch. The integration of the two pillars is performed by a deterministic rule represented as two successive decision tables that produce, first, the class of historical evidence of raw-material origin and then the final diagnosis among the classes of raw-material suspicion, equipment suspicion, parallel investigation, and insufficient evidence. The architecture prioritizes interpretability and auditability over \emph{end-to-end} optimization, in contrast to learned fusion models that would not be trainable in a scenario without standardized labels. A systematic literature review of Computing applied to industrial systems was conducted over the IEEE Xplore, ACM Digital Library, and Scopus databases, with a final selection of sixteen related works that evidence the gap addressed by the proposal. Retrospective validation over events with cause documented by expert inspection and deployment under a continuous inference regime, integrated with the shop-floor system, constitute the work planned between this qualification and the defense.

   \textbf{Keywords}: wire breakage. anomaly detection. historical statistics. evidence fusion. industrial decision support.
 \end{otherlanguage*}
\end{resumo}
